% GENERATED FILE. Edit canonical lessons, then run npm run generate:books.
% canonical-source-hash: fnv1a64:e680c4ee20e640fa
% canonical-lessons: JA-C221-taion, JA-C221-sagaru, JA-C221-itamu, JA-C221-chiryou, JA-C221-yobou, JA-R221-review-216-221

\chapter{Fever, Pain and Treatment}
\label{ch:ja-fever-pain-and-treatment}
\hlchaptermodality{221}

\begin{chapteropening}
\textbf{By the end of this chapter:} \emph{I can ask about a fever and say where it aches.}

Complete the last lesson of chapter 221: I can ask about a fever and say where it aches.
\end{chapteropening}

\section[\texorpdfstring{\ja{たいおん} (taion)}{たいおん (taion)}]{\ja{たいおん} (taion) — body temperature}

\label{lesson:JA-C221-taion}



\begin{quote}
Before the new one: say the Japanese for to lose weight, then the Japanese for an appetite.
\end{quote}

\subsection*{You'll want to know: \ja{たいおん}}
\textbf{\ja{たいおん}} — \emph{taion} — \textquotedblleft{}body temperature\textquotedblright{}.

\begin{quote}
\textbf{\ja{たいおんは} \ja{たかいですか。}} — \emph{taion wa takai desu ka} — \textquotedblleft{}Is your temperature high?\textquotedblright{}
\end{quote}

\begin{cousinweb}
\textbf{\ja{たいおん}} (\emph{taion}) is Sino-Japanese (kango): two Chinese characters read with their on'yomi, readings borrowed from Chinese. They are \emph{tai}, \textquotedblleft{}body\textquotedblright{} (Middle Chinese \emph{thejX}) and \emph{on}, \textquotedblleft{}warm\textquotedblright{} (\emph{'won}) (KANJIDIC2; Middle Chinese in Baxter's transcription).
\end{cousinweb}

\begin{grammarlens}[title={What you've built}]
One more word for everyday life at A2.
\end{grammarlens}

\subsection*{Your turn}
\begin{itemize}
\raggedright
  \item \emph{Say it:} \emph{taion}
  \item \emph{Say it:} \emph{taion}, once more
  \item \emph{Say it:} \emph{taion wa takai desu ka} — Is your temperature high?
  \item \emph{Recall:} say the Japanese for to lose weight, then the Japanese for an appetite, then say \emph{taion} again
\end{itemize}

\begin{tcolorbox}[breakable,colback=teal!4,colframe=teal!35!black,title={Before you move on}]
What does \ja{たいおん} mean? (\textquotedblleft{}Body temperature\textquotedblright{}.) Say it once more.
\end{tcolorbox}
\section[\texorpdfstring{\ja{さがる} (sagaru)}{さがる (sagaru)}]{\ja{さがる} (sagaru) — to go down; to drop}

\label{lesson:JA-C221-sagaru}



\begin{quote}
Before the new one: say the Japanese for an appetite, then the Japanese for body temperature.
\end{quote}

\subsection*{You'll want to know: \ja{さがる}}
\textbf{\ja{さがる}} — \emph{sagaru} — \textquotedblleft{}to go down; to drop\textquotedblright{}.

\begin{quote}
\textbf{\ja{ねつは} \ja{さがりましたか。}} — \emph{netsu wa sagarimashita ka} — \textquotedblleft{}Has your fever gone down?\textquotedblright{}
\end{quote}

Polite \textbf{\ja{さがります}}, \emph{sagarimasu}; a godan verb.

\begin{cousinweb}
\textbf{\ja{さがる}} (\emph{sagaru}) is a native Japanese word (wago), inherited rather than borrowed from Chinese; its kanji was attached later as a kun'yomi, and its older origin is not settled (JMdict, KANJIDIC2).
\end{cousinweb}

\begin{grammarlens}[title={What you've built}]
One more word for everyday life at A2.
\end{grammarlens}

\subsection*{Your turn}
\begin{itemize}
\raggedright
  \item \emph{Say it:} \emph{sagaru}
  \item \emph{Say it:} \emph{sagaru}, once more
  \item \emph{Say it:} \emph{netsu wa sagarimashita ka} — Has your fever gone down?
  \item \emph{Recall:} say the Japanese for an appetite, then the Japanese for body temperature, then say \emph{sagaru} again
\end{itemize}

\begin{tcolorbox}[breakable,colback=teal!4,colframe=teal!35!black,title={Before you move on}]
What does \ja{さがる} mean? (\textquotedblleft{}To go down; to drop\textquotedblright{}.) Say it once more.
\end{tcolorbox}
\section[\texorpdfstring{\ja{いたむ} (itamu)}{いたむ (itamu)}]{\ja{いたむ} (itamu) — to ache; to hurt}

\label{lesson:JA-C221-itamu}



\begin{quote}
Before the new one: say the Japanese for body temperature, then the Japanese for to go down.
\end{quote}

\subsection*{You'll want to know: \ja{いたむ}}
\textbf{\ja{いたむ}} — \emph{itamu} — \textquotedblleft{}to ache; to hurt\textquotedblright{}.

\begin{quote}
\textbf{\ja{はが} \ja{いたみますか。}} — \emph{ha ga itamimasu ka} — \textquotedblleft{}Does your tooth ache?\textquotedblright{}
\end{quote}

Polite \textbf{\ja{いたみます}}, \emph{itamimasu}; a godan verb. The noun \textbf{\ja{いたみ}}, \emph{itami}, means \textquotedblleft{}pain\textquotedblright{}.

\begin{cousinweb}
\textbf{\ja{いたむ}} (\emph{itamu}) is a native Japanese word (wago), inherited rather than borrowed from Chinese; its kanji was attached later as a kun'yomi, and its older origin is not settled (JMdict, KANJIDIC2).
\end{cousinweb}

\begin{grammarlens}[title={What you've built}]
One more word for everyday life at A2.
\end{grammarlens}

\subsection*{Your turn}
\begin{itemize}
\raggedright
  \item \emph{Say it:} \emph{itamu}
  \item \emph{Say it:} \emph{itamu}, once more
  \item \emph{Say it:} \emph{ha ga itamimasu ka} — Does your tooth ache?
  \item \emph{Recall:} say the Japanese for body temperature, then the Japanese for to go down, then say \emph{itamu} again
\end{itemize}

\begin{tcolorbox}[breakable,colback=teal!4,colframe=teal!35!black,title={Before you move on}]
What does \ja{いたむ} mean? (\textquotedblleft{}To ache; to hurt\textquotedblright{}.) Say it once more.
\end{tcolorbox}
\section[\texorpdfstring{\ja{ちりょう} (chiryō)}{ちりょう (chiryō)}]{\ja{ちりょう} (chiryō) — treatment; medical care}

\label{lesson:JA-C221-chiryou}



\begin{quote}
Before the new one: say the Japanese for to go down, then the Japanese for to ache.
\end{quote}

\subsection*{You'll want to know: \ja{ちりょう}}
\textbf{\ja{ちりょう}} — \emph{chiryō} — \textquotedblleft{}treatment; medical care\textquotedblright{}.

\begin{quote}
\textbf{\ja{ちりょうは} \ja{まだ} \ja{おわりません。}} — \emph{chiryō wa mada owarimasen} — \textquotedblleft{}The treatment is not over yet.\textquotedblright{}
\end{quote}

\begin{cousinweb}
\textbf{\ja{ちりょう}} (\emph{chiryō}) is Sino-Japanese (kango): two Chinese characters read with their on'yomi, readings borrowed from Chinese. They are \emph{chi}, \textquotedblleft{}to govern, cure\textquotedblright{} (Middle Chinese \emph{dri}) and \emph{ryō}, \textquotedblleft{}to heal\textquotedblright{} (\emph{ljewH}) (KANJIDIC2; Middle Chinese in Baxter's transcription).
\end{cousinweb}

\begin{grammarlens}[title={What you've built}]
One more word for everyday life at A2.
\end{grammarlens}

\subsection*{Your turn}
\begin{itemize}
\raggedright
  \item \emph{Say it:} \emph{chiryō}
  \item \emph{Say it:} \emph{chiryō}, once more
  \item \emph{Say it:} \emph{chiryō wa mada owarimasen} — The treatment is not over yet.
  \item \emph{Recall:} say the Japanese for to go down, then the Japanese for to ache, then say \emph{chiryō} again
\end{itemize}

\begin{tcolorbox}[breakable,colback=teal!4,colframe=teal!35!black,title={Before you move on}]
What does \ja{ちりょう} mean? (\textquotedblleft{}Treatment; medical care\textquotedblright{}.) Say it once more.
\end{tcolorbox}
\section[\texorpdfstring{\ja{よぼう} (yobō)}{よぼう (yobō)}]{\ja{よぼう} (yobō) — prevention}

\label{lesson:JA-C221-yobou}



\begin{quote}
Before the new one: say the Japanese for to ache, then the Japanese for treatment.
\end{quote}

\subsection*{You'll want to know: \ja{よぼう}}
\textbf{\ja{よぼう}} — \emph{yobō} — \textquotedblleft{}prevention\textquotedblright{}.

\begin{quote}
\textbf{\ja{よぼうの} \ja{ちゅうしゃですか。}} — \emph{yobō no chūsha desu ka} — \textquotedblleft{}Is it a preventive injection?\textquotedblright{}
\end{quote}

\begin{cousinweb}
\textbf{\ja{よぼう}} (\emph{yobō}) is Sino-Japanese (kango): two Chinese characters read with their on'yomi, readings borrowed from Chinese. They are \emph{yo}, \textquotedblleft{}beforehand\textquotedblright{} (Middle Chinese \emph{yoH}) and \emph{bō}, \textquotedblleft{}to ward off\textquotedblright{} (\emph{bjang}) (KANJIDIC2; Middle Chinese in Baxter's transcription).
\end{cousinweb}

\begin{grammarlens}[title={What you've built}]
One more word for everyday life at A2.
\end{grammarlens}

\subsection*{Your turn}
\begin{itemize}
\raggedright
  \item \emph{Say it:} \emph{yobō}
  \item \emph{Say it:} \emph{yobō}, once more
  \item \emph{Say it:} \emph{yobō no chūsha desu ka} — Is it a preventive injection?
  \item \emph{Recall:} say the Japanese for to ache, then the Japanese for treatment, then say \emph{yobō} again
\end{itemize}

\begin{tcolorbox}[breakable,colback=teal!4,colframe=teal!35!black,title={Before you move on}]
What does \ja{よぼう} mean? (\textquotedblleft{}Prevention\textquotedblright{}.) Say it once more.
\end{tcolorbox}
\section[(dialogue)]{Words from chapters 216–221 — a review}

\label{lesson:JA-R221-review-216-221}



\begin{quote}
Say the Japanese for an illness and condition.
\end{quote}

\subsection*{Your turn}
Say these aloud:

\begin{itemize}
\raggedright
  \item an illness, condition, to get tired, to get better, to shiver
  \item a patient, a consultation, an injection, a dentist, a cavity
  \item going into hospital, leaving hospital, surgery, a wound, life
  \item to fall over, to slip, to snap, to collapse, an injury
  \item health, robust, to gain weight, to lose weight, an appetite
  \item body temperature, to go down, to ache, treatment, prevention
\end{itemize}

\begin{tcolorbox}[breakable,colback=teal!4,colframe=teal!35!black,title={Before you move on}]
An illness? (\textbf{\ja{びょうき}}, \emph{byōki}.) Condition? (\textbf{\ja{ぐあい}}, \emph{guai}.) To get tired? (\textbf{\ja{つかれる}}, \emph{tsukareru}.) To get better? (\textbf{\ja{なおる}}, \emph{naoru}.) To shiver? (\textbf{\ja{ふるえる}}, \emph{furueru}.) A patient? (\textbf{\ja{かんじゃ}}, \emph{kanja}.) A consultation? (\textbf{\ja{しんさつ}}, \emph{shinsatsu}.) An injection? (\textbf{\ja{ちゅうしゃ}}, \emph{chūsha}.) A dentist? (\textbf{\ja{はいしゃ}}, \emph{haisha}.) A cavity? (\textbf{\ja{むしば}}, \emph{mushiba}.) Going into hospital? (\textbf{\ja{にゅういん}}, \emph{nyūin}.) Leaving hospital? (\textbf{\ja{たいいん}}, \emph{taiin}.) Surgery? (\textbf{\ja{しゅじゅつ}}, \emph{shujutsu}.) A wound? (\textbf{\ja{きず}}, \emph{kizu}.) Life? (\textbf{\ja{いのち}}, \emph{inochi}.) To fall over? (\textbf{\ja{ころぶ}}, \emph{korobu}.) To slip? (\textbf{\ja{すべる}}, \emph{suberu}.) To snap? (\textbf{\ja{おれる}}, \emph{oreru}.) To collapse? (\textbf{\ja{たおれる}}, \emph{taoreru}.) An injury? (\textbf{\ja{けが}}, \emph{kega}.) Health? (\textbf{\ja{けんこう}}, \emph{kenkō}.) Robust? (\textbf{\ja{じょうぶ}}, \emph{jōbu}.) To gain weight? (\textbf{\ja{ふとる}}, \emph{futoru}.) To lose weight? (\textbf{\ja{やせる}}, \emph{yaseru}.) An appetite? (\textbf{\ja{しょくよく}}, \emph{shokuyoku}.) Body temperature? (\textbf{\ja{たいおん}}, \emph{taion}.) To go down? (\textbf{\ja{さがる}}, \emph{sagaru}.) To ache? (\textbf{\ja{いたむ}}, \emph{itamu}.) Treatment? (\textbf{\ja{ちりょう}}, \emph{chiryō}.) Prevention? (\textbf{\ja{よぼう}}, \emph{yobō}.)
\end{tcolorbox}
